\documentclass[12pt,a4paper]{article}
\usepackage[T1]{fontenc}
\usepackage[utf8]{inputenc}
\usepackage{mathptmx}
\renewcommand{\figurename}{Figura}
\renewcommand{\tablename}{Tabela}
\usepackage[a4paper,left=3cm,right=3cm,top=2.5cm,bottom=2.5cm]{geometry}
\usepackage{amsmath,amssymb,amsthm}
\usepackage{graphicx}
\usepackage{booktabs}
\usepackage{algorithm}
\usepackage{algpseudocode}
\usepackage{tikz}
\usepackage[square,authoryear]{natbib}
\usepackage[colorlinks=true,linkcolor=black,citecolor=black,urlcolor=black]{hyperref}
\usepackage{url}
\usepackage{caption}
\captionsetup{font=small,labelfont=bf,labelsep=period,justification=centering}
\usetikzlibrary{positioning,arrows.meta}

\setlength{\parindent}{1.0cm}
\setlength{\parskip}{0pt}
\renewcommand{\baselinestretch}{1.0}

% títulos de seção no estilo SBC (negrito, sem serifa especial)
\makeatletter
\renewcommand{\section}{\@startsection{section}{1}{0pt}{-14pt plus -2pt minus -1pt}{6pt}{\normalfont\large\bfseries}}
\renewcommand{\subsection}{\@startsection{subsection}{2}{0pt}{-10pt plus -2pt minus -1pt}{4pt}{\normalfont\normalsize\bfseries}}
\makeatother

% resumo e abstract no estilo SBC (itálico, recuado)
\newenvironment{resumobloco}[1]
  {\begin{list}{}{\setlength{\leftmargin}{0.8cm}\setlength{\rightmargin}{0.8cm}\setlength{\itemindent}{0pt}\setlength{\listparindent}{0pt}}
   \item[]\small\itshape\textbf{#1.}\ }
  {\end{list}}

\theoremstyle{plain}
\newtheorem{lema}{Lema}
\newtheorem{teorema}{Teorema}
\theoremstyle{definition}
\newtheorem*{prova}{Demonstração}

% citações no formato do Trabalho 1: [Autor et al. 2002]
\renewcommand{\cite}{\citep}
\bibpunct[, ]{[}{]}{;}{a}{}{,}
\algrenewcommand\algorithmicrequire{\textbf{Entrada:}}
\algrenewcommand\algorithmicensure{\textbf{Saída:}}
\algrenewcommand\algorithmicif{\textbf{se}}
\algrenewcommand\algorithmicthen{\textbf{então}}
\algrenewcommand\algorithmicfor{\textbf{para}}
\algrenewcommand\algorithmicdo{\textbf{faça}}
\algrenewcommand\algorithmicreturn{\textbf{devolva}}
\algrenewcommand\algorithmicwhile{\textbf{enquanto}}
\algrenewcommand\algorithmicelse{\textbf{senão}}
\algrenewcommand\algorithmicend{\textbf{fim}}
\algrenewcommand\algorithmicforall{\textbf{para todo}}
\algrenewtext{EndIf}{\algorithmicend\ \algorithmicif}
\algrenewtext{EndFor}{\algorithmicend\ \algorithmicfor}
\algrenewtext{EndWhile}{\algorithmicend\ \algorithmicwhile}
\algrenewtext{EndProcedure}{\algorithmicend\ \algorithmicprocedure}
\renewcommand{\algorithmicprocedure}{\textbf{procedimento}}
\floatname{algorithm}{Algoritmo}
